\label{sec:theorem}

The empirical claim of this paper is that a single LC-PINN
$u_\theta(x,\lambda)$, trained with $\lambda$ sampled from a probability
measure $p_\lambda$, parameterises the parametric PDE family
$\{F_\lambda u = 0\}_{\lambda \in \Lambda}$ \emph{simultaneously}
in~$\lambda$. The following proposition makes precise the sense in
which this holds at a global optimum of the LC-PINN objective.

\paragraph{Setup.}
Let $\Lambda \subseteq \mathbb{R}^{d_\lambda}$ be the parameter set,
$p_\lambda$ a probability measure on $\Lambda$ with full support, and
$\Omega \subseteq \mathbb{R}^{d_x}$ the spatial--temporal domain with
sampling measure $p_\Omega$. Each $\lambda \in \Lambda$ indexes a
residual operator $F_\lambda$ together with associated boundary,
initial, and (optional) data-fit operators
$\{B_\lambda, I_\lambda, D_\lambda\}$. For a candidate field
$v: \Omega \to \mathbb{R}$, write the per-$\lambda$ functional
\[
   \mathcal{L}(\lambda; v) \;=\;
     w_F\,\| F_\lambda v \|_{L^2(\Omega, p_\Omega)}^2
   + w_B\,\| B_\lambda v \|_{L^2(\partial\Omega)}^2
   + w_I\,\| I_\lambda v \|_{L^2(t=0)}^2
   + w_D\,\| D_\lambda v \|_{L^2}^2 ,
\]
with strictly positive loss weights $(w_F, w_B, w_I, w_D)$. The
LC-PINN training objective is the $p_\lambda$-expectation
\[
   J(\theta) \;=\;
     \int_{\Lambda} \mathcal{L}\bigl(\lambda;\, u_\theta(\,\cdot\,, \lambda)\bigr)
     \, dp_\lambda(\lambda) .
\]
Assume:
\begin{itemize}\setlength{\itemsep}{0pt}
   \item[(A1)] the network class
         $\{u_\theta(\cdot, \lambda) : \theta \in \Theta\}$ contains,
         for every $\lambda$, a residual-minimiser
         $u^\star_\lambda \in \arg\min_v \mathcal{L}(\lambda; v)$ with
         $\mathcal{L}(\lambda; u^\star_\lambda) = 0$;
   \item[(A2)] the map
         $\lambda \mapsto \mathcal{L}(\lambda; u_\theta(\cdot,\lambda))$
         is $dp_\lambda$-integrable for every $\theta$;
   \item[(A3)] $u_\theta$ is jointly continuous in $(x, \lambda)$.
\end{itemize}

\begin{proposition}[Residual-minimiser at the LC optimum]
\label{prop:lc-invariance}
Let $\theta^\star \in \arg\min_\theta J(\theta)$. Under
\textnormal{(A1)--(A3)}, for $p_\lambda$-almost every
$\lambda \in \Lambda$ the section
$u_{\theta^\star}(\cdot, \lambda)$ is a residual-minimiser of the
parametric PDE indexed by~$\lambda$:
\[
   \mathcal{L}\bigl(\lambda;\, u_{\theta^\star}(\,\cdot\,, \lambda)\bigr)
   \;=\; 0
   \qquad\text{for }\; p_\lambda\text{-a.e.\ } \lambda \in \Lambda .
\]
In particular, on $\operatorname{supp}(p_\lambda)$ the LC-PINN section
satisfies $F_\lambda u_{\theta^\star}(\cdot,\lambda) = 0$ in
$L^2(\Omega, p_\Omega)$ and the auxiliary constraints to the same
precision.
\end{proposition}

\paragraph{Sketch proof.}
The integrand
$\lambda \mapsto \mathcal{L}(\lambda; u_\theta(\cdot,\lambda))$ is
non-negative for every $\theta$ and every $\lambda$, since each summand
is a squared $L^2$-norm. By (A1) the per-$\lambda$ minimum is zero, so
$\inf_\theta J(\theta) \ge 0$ and the lower bound is attainable: any
measurable selector $\lambda \mapsto \theta_\lambda$ with
$u_{\theta_\lambda}(\cdot, \lambda) = u^\star_\lambda$ exists pointwise
under (A1), and a single $\theta^\star$ achieving $J(\theta^\star) = 0$
exists when the network class is rich enough to parameterise the family
$\{u^\star_\lambda\}_{\lambda \in \Lambda}$ jointly (e.g.\ universal
approximation in $C(\Omega \times \Lambda)$). For such $\theta^\star$,
\[
   0 \;=\; J(\theta^\star)
   \;=\; \int_{\Lambda}
            \mathcal{L}\bigl(\lambda;\, u_{\theta^\star}(\,\cdot\,, \lambda)\bigr)
         \, dp_\lambda(\lambda) ,
\]
and the integrand is non-negative; the standard
``vanishing-integral of a non-negative function'' lemma then gives
$\mathcal{L}(\lambda; u_{\theta^\star}(\cdot,\lambda)) = 0$ for
$p_\lambda$-a.e.~$\lambda$. \hfill$\square$

\paragraph{Remark 1 (where the support assumption matters).}
The conclusion is a $p_\lambda$-a.e.\ statement on $\Lambda$: nothing
prevents the LC-PINN from being arbitrarily wrong on a
$p_\lambda$-null set, including the boundary of $\Lambda$ if
$p_\lambda$ has no mass there. In our experiments
$p_\lambda = \mathrm{Uniform}(\Lambda)$ has full support, so the
$p_\lambda$-a.e.\ statement coincides with Lebesgue-a.e.\ on
$\Lambda$.

\begin{corollary}[Finite-capacity LC bound]
\label{cor:finite-capacity}
Replace assumption \textnormal{(A1)} by the weaker
\emph{$\varepsilon$-residual capacity} condition: there exists
$\bar\theta \in \Theta$ with
$\mathcal{L}(\lambda; u_{\bar\theta}(\cdot,\lambda)) \le \varepsilon$
for $p_\lambda$-a.e.\ $\lambda$. Suppose the optimiser reaches an
objective value within $\delta$ of the infimum:
$J(\theta^\star) \le \inf_\theta J(\theta) + \delta$. Then for every
$\eta > 0$,
\[
   \Pr_{\lambda \sim p_\lambda}\!\left[
      \mathcal{L}\bigl(\lambda;\, u_{\theta^\star}(\,\cdot\,, \lambda)\bigr)
      \;>\; \varepsilon + \eta
   \right]
   \;\le\; \frac{\delta}{\eta} ,
\]
i.e.\ on a $p_\lambda$-set of mass $\ge 1 - \delta/\eta$ the
per-$\lambda$ residual loss is at most $\varepsilon + \eta$.
\end{corollary}

\paragraph{Sketch.}
By definition of $\bar\theta$ and (A2),
$\inf_\theta J(\theta) \le J(\bar\theta) \le \varepsilon$, so
$J(\theta^\star) \le \varepsilon + \delta$. Let
$g(\lambda) = \mathcal{L}(\lambda; u_{\theta^\star}(\cdot,\lambda)) -
\varepsilon$, which is $p_\lambda$-a.e.\ bounded below by
$-\varepsilon$ but in particular below by zero on the
$p_\lambda$-set where the comparator $\bar\theta$ achieves the
$\varepsilon$-bound. Markov's inequality applied to the
non-negative part $g_+ = \max(g, 0)$ gives
$\Pr[g_+ > \eta] \le \mathbb{E}[g_+]/\eta \le \delta/\eta$.
\hfill$\square$

This is the practically observable form of the
$\lambda$-invariance result: it bounds the per-$\lambda$ loss for
\emph{any} suboptimal $\theta^\star$ by the realisable residual
budget plus the optimisation gap, and matches the per-$\lambda$
rel-$L^2$ tables in \Cref{sec:results}.

\paragraph{Remark 2 (sampling-distribution invariance).}
If $p_\lambda$ and $\tilde p_\lambda$ are two probability measures
with the same support $\Lambda_0 \subseteq \Lambda$, the sets of
\emph{a.e.-residual-minimisers} on $\Lambda_0$ coincide. So the
optimum-set of the LC objective on $\operatorname{supp}(p_\lambda)$
does not depend on the precise sampling law---only on the support.
This is an asymptotic statement; at finite training budget the law
of $p_\lambda$ governs which slices receive enough gradient signal,
which we observe empirically (\Cref{sec:results-ablations}).

\paragraph{Remark 3 (sample complexity per step).}
\label{rem:sample}
At each Adam step the gradient is built from $n_\lambda$ i.i.d.\
draws of $\lambda \sim p_\lambda$. Conditional on $\theta$, the
per-step gradient estimator
$\hat g_{n_\lambda} = \frac{1}{n_\lambda}\sum_{i=1}^{n_\lambda}
\nabla_\theta \mathcal{L}(\lambda_i; u_\theta(\cdot,\lambda_i))$
is unbiased for $\nabla_\theta J(\theta)$, and if the per-$\lambda$
gradient norm is $p_\lambda$-essentially bounded by $G$ then
$\Pr[\| \hat g_{n_\lambda} - \nabla J \|_\infty > t] \le
2 \exp(-n_\lambda t^2 / 2 G^2)$ by Hoeffding. The variance of the
estimate is $O(1/n_\lambda)$, so the regime $n_\lambda \!\sim\! 4$
empirically chosen in \Cref{sec:results-ablations} corresponds to
the sweet spot between gradient noise (small $n_\lambda$) and
budget concentration (large $n_\lambda$ at fixed total step count).
